\ticket{22}{Дистрибутивная семантика, Word2Vec и FastText}

\begin{examplan}
    \item Объяснить дистрибутивную гипотезу и понятие эмбеддинга.
    \item Описать Word2Vec: CBOW, Skip-gram и способы обучения.
    \item Описать FastText и учет подслов.
    \item Раскрыть понятие семантического пространства слов.
\end{examplan}

\subsection*{Короткая версия для письменного ответа}

\textbf{Дистрибутивная семантика} основана на гипотезе: слова, встречающиеся в похожих контекстах, имеют похожие значения. Поэтому значение слова можно представить числовым вектором, обученным на контекстах его употребления.

\textbf{Эмбеддинг} --- плотный вектор слова в многомерном семантическом пространстве.

\subsection*{Свойства эмбеддингов}

\begin{itemize}
    \item близкие по смыслу слова имеют близкие векторы;
    \item векторы можно использовать как признаки в NLP-задачах;
    \item некоторые отношения отражаются как направления в пространстве;
    \item близость часто измеряют косинусной мерой.
\end{itemize}

Классический пример аналогии: \textit{king} $-$ \textit{man} $+$ \textit{woman} $\approx$ \textit{queen}.

\subsection*{Word2Vec}

\textbf{Word2Vec} --- семейство нейросетевых моделей для обучения статических эмбеддингов слов. Основные архитектуры: CBOW и Skip-gram.

\textbf{CBOW} предсказывает центральное слово по окружающим словам. Контекстные слова переводятся в эмбеддинги, усредняются, затем модель предсказывает слово. CBOW обычно быстрее и хорошо работает с частыми словами.

\textbf{Skip-gram} предсказывает контекстные слова по центральному слову. Он часто лучше работает с редкими словами, но обучается медленнее.

Для ускорения обучения вместо полного softmax применяются:
\begin{itemize}
    \item \textbf{hierarchical softmax}: слово кодируется путем в дереве;
    \item \textbf{negative sampling}: модель отличает реальные пары слово-контекст от случайных отрицательных пар.
\end{itemize}

\subsection*{FastText}

\textbf{FastText} расширяет Word2Vec: слово представляется как набор символьных n-грамм. Вектор слова получается из векторов его подсловных частей.

Преимущества:
\begin{itemize}
    \item лучше работает с редкими словами;
    \item может строить векторы для слов вне словаря;
    \item учитывает морфологию, что важно для языков с богатым словоизменением;
    \item полезен для ошибок, новых слов и производных форм.
\end{itemize}

FastText может обучаться в вариантах CBOW или Skip-gram, но вместо только целых слов использует также подсловные n-граммы.

\subsection*{Семантическое пространство}

Обученные векторы образуют семантическое пространство, где:
\begin{itemize}
    \item близость точек соответствует сходству значений;
    \item направления могут соответствовать отношениям: род, число, пол, страна-столица;
    \item пространство можно визуализировать через PCA или t-SNE.
\end{itemize}

Word2Vec и FastText дают \textbf{статические} эмбеддинги: у слова один вектор независимо от контекста. Современные модели вроде BERT строят \textbf{контекстуальные} эмбеддинги, где вектор слова зависит от предложения.

\begin{important}
    \item Дистрибутивная гипотеза: значение слова определяется контекстами.
    \item Эмбеддинг --- плотный числовой вектор слова.
    \item CBOW предсказывает слово по контексту, Skip-gram --- контекст по слову.
    \item FastText использует символьные n-граммы и лучше справляется с редкими/OOV словами.
    \item Статические эмбеддинги не различают значения многозначного слова по контексту.
\end{important}

